%%
%% Copyright 2022 OXFORD UNIVERSITY PRESS
%%
%% This file is part of the 'oup-authoring-template Bundle'.
%% ---------------------------------------------
%%
%% It may be distributed under the conditions of the LaTeX Project Public
%% License, either version 1.2 of this license or (at your option) any
%% later version.  The latest version of this license is in
%%    http://www.latex-project.org/lppl.txt
%% and version 1.2 or later is part of all distributions of LaTeX
%% version 1999/12/01 or later.
%%
%% The list of all files belonging to the 'oup-authoring-template Bundle' is
%% given in the file `manifest.txt'.
%%
%% Template article for OXFORD UNIVERSITY PRESS's document class `oup-authoring-template'
%% with bibliographic references
%%


\documentclass[unnumsec,webpdf,namedate,modern,large]{oup-authoring-template}%

% Hyperlinks
\usepackage{hyperref}
\hypersetup{
    colorlinks=true,
    urlcolor=blue,
    citecolor=black,
    linkcolor=blue,
    }

\usepackage{graphicx}
\setcounter{figure}{0}
\graphicspath{{figures/}}

% New command to hyperlink figures
% See https://tex.stackexchange.com/questions/266434/extend-the-hyperref-link-to-figure-and-a-b-c
\newcommand*{\figref}[2][]{%
  \hyperref[{fig:#2}]{%
    Fig~\ref*{fig:#2}%
    \ifx\\#1\\%
    \else
      \,#1%
    \fi
  }%
}

% Reduce vertical spacing between sections
\usepackage{titlesec}
\titlespacing*{\section}{0pt}{0.2\baselineskip}{0.2\baselineskip}

%%% Bibliography %%%
\usepackage[style=authoryear,backend=biber,maxnames=1,uniquelist=false]{biblatex}
\addbibresource{references.bib}
                           
%% Bibiography style for articles %%%
\DeclareFieldFormat[article]{title}{#1} % Remove quotes from article titles
\DeclareFieldFormat[article]{volume}{\mkbibbold{#1}} % Make volume bold
\DeclareFieldFormat[article]{pages}{#1} % Remove "pp" abbreviation before the page number

% Customize bibliography according to Mini Oxford SCIMED style for Bioinformatics journal
\DeclareBibliographyDriver{article}{%
  \printnames{author}%%
  \newunit
  \printfield{title}.%
  \newunit
  \printfield{journaltitle}%
  \setunit*{\space}%
  \printfield{year}%
  \iffieldundef{volume}
    {\adddot} % If volume is undefined, end with a dot
    {%
      \addsemicolon\printfield{volume} % Add a semicolon, print volume
      \iffieldundef{pages}
        {\adddot} % If pages are undefined, end with a dot
         {\addcolon\printfield{pages}\adddot} % If pages are defined, print colon, pages, and end with a dot
     }
 }

%% Bibiography style for manuals %%
\DeclareFieldFormat[manual]{volume}{\mkbibbold{#1}} % Make volume bold
\DeclareFieldFormat[manual]{pages}{#1} % Remove "pp" abbreviation before the page number

% For beautiful tables
\usepackage{booktabs}

% Remove white space above and below the tables
\usepackage{float}
\setlength{\floatsep}{10pt}
\setlength{\textfloatsep}{10pt}
\setlength{\intextsep}{10pt}

% Reduce blank space between sections
\usepackage[skip=1pt]{parskip}

\begin{document}

\journaltitle{Bioinformatics}

\title{HoloFoodR: a statistical programming framework for holo-omics data integration workflows}

\firstpage{1}

\author[1,\dag]{Tuomas Borman\ORCID{0000-0002-8563-8884}}
\author[2,\dag]{Artur Sannikov\ORCID{0000-0001-7765-123X}}
\author[3]{Robert D. Finn\ORCID{0000-0001-8626-2148}}
\author[4]{Morten Tønsberg Limborg\ORCID{0000-0002-7718-6531}}
\author[5]{Alexander B. Rogers\ORCID{0000-0002-4283-6135}}
\author[6]{Varsha Kale\ORCID{0000-0002-3503-6383}}
\author[7]{Kati Hanhineva\ORCID{0000-0001-6834-7375}}
\author[8]{Leo Lahti\ORCID{0000-0001-5537-637X}}

\address[1,8]{\orgdiv{Department of Computing}, \orgname{University of Turku}, \postcode{20014} \state{Turku}, \country{Finland}}
\address[2,7]{\orgdiv{Department of Life Technologies}, \orgname{University of Turku}, \postcode{20014} \state{Turku}, \country{Finland}}
\address[3,5,6]{\orgdiv{Wellcome Genome Campus}, \orgname{EMBL-EBI}, \postcode{CB10 1SA} \state{Hinxton, Cambridgeshire}, \country{United Kingdom}}
\address[4]{\orgdiv{Center for Evolutionary Hologenomics, GLOBE Institute, Faculty of Health and Medical Sciences}, \orgname{University of Copenhagen}, \postcode{1353} \state{Copenhagen}, \country{Denmark}}
\corresp{\dag = Equal contribution.}

% Abstract
\abstract{
\textbf{Summary:} Holo-omics is an emerging research area that integrates multi-omic datasets from the host organism and its microbiome to study their interactions. Recently, curated and openly accessible holo-omic databases have been developed. The HoloFood database, for instance, provides nearly 10,000 holo-omic profiles for salmon and chicken under controlled treatments. However, bridging the gap between holo-omic data resources and algorithmic frameworks remains a challenge. Combining the latest advances in statistical programming with curated holo-omic data sets can facilitate the design of open and reproducible research workflows in the emerging field of holo-omics.
\\ % New line
\textbf{Availability and implementation:} HoloFoodR R/Bioconductor package and the source code are available under the open-source Artistic License 2.0 at the package homepage \url{https://doi.org/10.18129/B9.bioc.HoloFoodR}.
\\ % New line
\textbf{Contact:} tuomas.v.borman@utu.fi
\\ % New line
\textbf{Supplementary information:} Supplementary information is available in the package vignette \url{https://ebi-metagenomics.github.io/HoloFoodR/articles/case_study.html}
}

\keywords{holo-omics, multi-omics, metagenomics, metabolomics, data integration, bioconductor}

\maketitle

\section{Introduction}

The rapid advancement of omics technologies, including (meta)genomics and metabolomics, has been driven by breakthroughs in computational methods (\cite{moreno-indias_statistical_2021,marcos-zambrano_toolbox_2023,santamaria_bioinformatic_2024}).
These developments have enabled holo-omics, a rapidly emerging field that uses an integrative approach to comprehensively collect and analyse omic data from both a host organism and its associated microbiome, collectively referred to as the {\it holobiont} (\cite{nyholmHoloOmicsIntegratedHostMicrobiota2020,limborgAppliedHologenomicsFeasibility2018,odriozola_practical_2024}). The holo-omic approach has improved our understanding of complex biological systems, for example, in aquaculture (\cite{limborgAppliedHologenomicsFeasibility2018,atlantic_salmon}).

Holo-omic research and the collaborative development of computational methods could benefit from curated, open-access data resources.
The limited availability of comprehensive multi-omic data resources can impede research progress, constraining computational method development. Open data portals with an accessible Application Programming Interface (API) are crucial in overcoming these challenges. They facilitate research by providing access to data resources and supporting the collaborative development and benchmarking of new data science methods to extract insights from large, curated datasets (\cite{pasolli_2017}).

Data can be limited in value unless it adheres to the FAIR principles (Findability, Accessibility, Interoperability, Reuse) (\cite{wilkinsonFAIRGuidingPrinciples2016}). API interfacing and data wrangling require specialized skills, leading to non-transferable, error-prone workflows not easily reusable by the wider scientific community. Therefore, standardised open-source workflows are needed to search, retrieve, and convert data into a suitable format for downstream analyses.
To narrow the gap between upstream data retrieval and downstream analysis, we developed the HoloFoodR package, which facilitates seamless programmatic linking of the holo-omic data via the HoloFood data portal API and analysis methods from the Bioconductor  (\cite{gentleman_bioconductor_2004,callahan_bioconductor_2016}).

\section{Materials and methods}

The HoloFood data portal is an open access portal of curated holo-omic data and analyses, developed by the international HoloFood consortium (\cite{rogers_holofood_2025}) and hosted by the European Bioinformatics Institute (EMBL-EBI). It centralizes access to heterogeneous data resources including the European Nucleotide Archive for sequence data, MGnify for metagenomic data, and MetaboLights for metabolomic data, and tracks their interrelations via a web portal and free API, covering nearly 10,000 samples from over 2,000 individual chickens and salmon. Biomolecular and physiological measurements were collected at the level of each individual specimen in the project to explore the effects of novel sustainable feeds on physiological processes in farmed animals. In addition, the metadata of each sample is stored in the EMBL-EBI BioSamples service (\cite{courtot_2018}).

We developed a conceptual framework to support best practices of statistical programming in holo-omic research, and implemented the work as the HoloFoodR R software library. The package relies on specialized data containers that have been designed to organize the multi-omic data in a structured format. Once the data has been imported into these formats, users can leverage state-of-the-art statistical programming methods for data processing and analysis.

\vspace{-1em}
\subsection{Integrative data containers}

Data containers provide structured, standardised storage for data, proving useful in life science informatics where complex, hierarchical and multi-source data collections are common to describe the studied phenomenon (\cite{drnevich_2025}). Custom data containers simplify handling diverse omics data, supporting the design of modular and reproducible workflows.

Two data structures are central to our approach, encompassing the needs to organize data for a single omics type, and subsequently link these across multiple omics.
First, the TreeSummarizedExperiment (TreeSE) data container (\cite{huang_treesummarizedexperiment_2021}) standardises single-omic data, such as metagenome or metabolome derived data sets. It stores experimental data, along with feature and sample metadata, in a structured format. Additionally, it supports the integration of hierarchical information on the features and samples, such as phylogenetic trees. The second container, MultiAssayExperiment (MAE) (\cite{ramos_software_2017}) acts as a logical complement by integrating multiple, heterogeneous omics datasets, seamlessly linking samples across single-omic datasets. Together, these data containers form the basis for an ecosystem of interoperable methods developed by the Bioconductor community.

\vspace{-1em}
\subsection{HoloFoodR workflow}

The HoloFoodR package structures data in these single and multi-omic formats, enabling efficient data importing from holo-omic databases for streamlined analysis and integration. The HoloFoodR workflow is summarized in \figref{workflow}. The package uses six functions to query and retrieve data from the HoloFood and MetaboLights databases (Table \ref{tab:holofoodr_functions}). A key feature is its ability to organize versatile data combinations into the MAE data container. Thus, it transforms complex data from the HoloFood API into a standardised structure for downstream statistical analyses in the Bioconductor data science environment  (\cite{ramos_software_2017}).

Furthermore, HoloFoodR leverages the MetaboLights database to retrieve metabolomic datasets indicated by project identifiers in the HoloFood data. The data can then be integrated with complementary metagenomic data resources that are available via the MGnify database using the MGnifyR package (\cite{mgnifyr}).

\vspace{-1em}

\begin{table}[H]
\begin{center}
\begin{minipage}{\linewidth}
\caption{HoloFoodR functions that facilitate data search and retrieval}\label{tab:holofoodr_functions}
    \begin{tabular}{l p{0.6\linewidth}}
        \toprule
        Function & Description \\ 
        \midrule
        addMGnify() & Integrate the results retrieved with getResult() with the metagenomic datasets fetched using the MGnifyR package. The result is in MultiAssayExperiment format. \\
        doQuery() & Search the HoloFood database for animals, genome catalogues, samples, or viral catalogues. \\
        getData() & Retrieve diverse data types from HoloFood database. Returns a list or data.frame, depending on the query. Offers more flexibility than getResult(), but returns unstandardised data structures. \\
        getMetaboLights() & Retrieve processed metabolomic data from MetaboLights as list or TreeSummarizedExperiment format. \\
        getMetaboLightsFile() & Downloads raw metabolomic data files from MetaboLights. \\
        getResult()  & Retrieve sample-level data (e.g., metadata and measurements) from the HoloFood database in MultiAssayExperiment format. \\
        \bottomrule
    \end{tabular}
\end{minipage}
\end{center}
\end{table}

\vspace{-1em}

\begin{figure*}[t]
	\centering
    % \vspace{20pt}%
    \includegraphics[width=1\textwidth]{holofoodr_workflow_9.png}
    \label{fig:workflow}
    \caption{Data science workflow for holo-omic analysis. Programmatic access to open omic databases is facilitated by the dedicated software that feeds the data to R computational environment. The Bioconductor packages are utilised to visualise microbial genera abundances and to analyse  time effect on omega-3 fatty acid concentrations in salmon muscles. The illustration and logo of the multi-assay data container have been adopted from the MultiAssayExperiment vignette (Artistic License 2.0).}
	\label{fig:workflow}
\end{figure*}

\vspace{-1em}
\subsection{Bioconductor methods ecosystem}

When data from the HoloFood database is retrieved via the HoloFoodR package, users gain direct access to single- and multi-omic downstream analysis methods available in Bioconductor. It currently offers 2,300 R packages, developed by a diverse community of individual contributors for a wide range of fields (see e.g. \cite{gentleman_bioconductor_2004,callahan_bioconductor_2016,amezquita_orchestrating_2020,drnevich_2025}). Bioconductor's \texttt{mia} framework centers around these containers and supports a range of analysis and visualisation methods for microbiome research (\cite{mia, miaViz, oma}). Similarly, the metabolomic data can be pre-processed, analysed and visualised using the \texttt{notame} package (\cite{klavus_2020}).

\section{Results}

This section presents a practical use case of applying HoloFoodR in a holo-omics workflow.
The workflow includes the following steps:

\begin{enumerate}
    \item Fetch and integrate data from the HoloFood and MGnify databases.
    \item Filter, clean, and transform data for analysis.
    \item Explore and summarize the data.
    \item Test associations between fatty acids, time, and treatment.
    \item Test associations between microbiome composition, time, and treatment.
    \item Characterise the joint variation between the parallel omics measurements.
\end{enumerate}

The full workflow, including figures and data refinement, is available in the package vignette at \url{https://ebi-metagenomics.github.io/HoloFoodR/articles/case_study.html}.

\vspace{-1em}
\subsection{Data import and wrangling}

The workflow begins by querying the HoloFood database for available animals. We gathered all salmon entries and their metadata to explore data types and samples. These samples were compiled into an MAE data container. For metagenomic data, we leveraged the capabilities of the MGnifyR package to identify the sample IDs from the MGnify database and fetch the metagenomic data as a MAE object. The consistent data format supports  seamless merging of HoloFood and MGnify datasets. The resulting container holds various types of omics with associated metadata for downstream analyses.

Using the MAE object for data wrangling, we focused on salmon from Trial A in the HoloFood database to study the effects of fermented seaweed as a feed additive. Samples were obtained from animals euthanised at either the start or end of the trial. Fatty acid concentrations were measured in muscle tissue, while metagenomic samples were taken from the intestine.

The \texttt{mia} framework (\cite{mia, miaViz}) provides an intuitive R interface designed for the downstream analysis, especially for microbiome data. We filtered and agglomerated the data to focus on specific microbes and fatty acids, applying centred log-ratio transformation to metagenomic data and logarithmic transformation to fatty acid data. The full details are provided in the HoloFoodR package vignette.

\vspace{-1em}
\subsection{Downstream analysis}

After data retrieval and basic exploration, we examined the effects of seaweed treatment and ageing. No treatment impact on fatty acid composition was observed, consistent with the recent HoloFood study (\cite{rasmussen_holo-omics_2025}). However, we found temporal effects: as the salmon grew, concentrations of several fatty acids (eicosapentaenoic, docosahexaenoic, linoleic, oleic, palmitic, and stearic) increased (see \figref{workflow}, data analysis panel, which highlights the first two fatty acids).

Microbial community analysis indicated a dominance of {\it Mycoplasma}, consistent with previous studies ({\cite{zarkasiPyrosequencingbasedCharacterizationGastrointestinal2014,bozzi_salmon_2021}}, see \figref{workflow}, data analysis panel). We also observed an increase in microbial diversity (Shannon index) with time, confirmed by Principal Coordinate Analysis (Bray-Curtis dissimilarity).

To explore the co-abundance of fatty acids and microbial genera, we performed multi-omic factor analysis (\cite{argelaguet_mofa_2020}), revealing that {\it Cetobacterium}, {\it Vibrio}, {\it Aliivibrio}, and {\it Photobacterium} covaried with overall fatty acid levels, while {\it Mycoplasma} showed no such correlation.

\section{Discussion and conclusions}
% Context / background / reflection on objectives

Integrating multiple omics layers provides holistic insights into complex systems. Holo-omics is a new field that emphasises the interactions between the hosts and their microbiomes. However, such analyses rely on curated data, computational environments, and algorithmic tools. Collaborative development and sharing in holo-omics can advance the field and prevent redundant efforts (\cite{open_microbiome, lahti_open_2018}).

To fill the gap between holo-omic data and downstream analysis methods, we developed the HoloFoodR package. Our work offers programmatic access to the HoloFood data portal and integrates with complementary data sources and Bioconductor analysis tools, supporting open data science in holo-omics.

HoloFoodR simplifies data retrieval by minimizing external data wrangling and adhering to standardised data structures that support efficient use of the multi-omic analyses methods. This standardisation replaces lengthy \textit{ad-hoc} code with concise and tested open-source solutions. We demonstrate this approach with an end-to-end workflow, highlighting the use of multi-assay data containers to facilitate access to Bioconductor methods.

% Use in teaching
Beyond method development and analysis, HoloFood data can aid in teaching multi-omics techniques. HoloFoodR offers simplified access to real-world omic data, enabling more proficient users to “learn by doing”. The complexity of this data makes it superior to “toy” datasets for teaching advanced data cleaning and analysis skills. (\cite{drnevich_2025}).

The HoloFood data portal currently provides the interconnectivity between the samples and omics datasets, based on BioSamples identifiers. This linkage provides a template for future multi-omics datasets, and as the number of datasets increases, the HoloFoodR package could be generalised so that the concepts can be readily applied to other datasets.

% Limitations, implementation issues
Despite the benefits, there are also limitations. First, raw spectral metabolite data from the MetaboLights database requires extensive preprocessing, often with external tools (\cite{klavus_2020}). Whereas the proposed custom data structures can readily support downstream analyses, their construction and use require sufficient R programming skills.

% Concluding sentences
Thus, the methods and open data science strategy that we have proposed can serve as a template for conducting multi-omic analyses. The HoloFoodR package can be adapted for other data resources relevant to holo-omic research, where the growing Bioconductor ecosystem offers an expanding set of data analysis tools. 

\section{Acknowledgements}

We are grateful to the HoloFood consortium (\url{https://www.holofood.eu}) for providing the curated database and API.

% \section{Supplementary data}

\section{Competing interests}
No competing interest is declared.

\section{Funding}
This work was supported by the European Commission in the framework of the Horizon2020 Project FindingPheno [GA  952914] and HoloFood [GA 817729]. A.S. and K.H. were  supported by Jane and Aatos Erkko Foundation and the Research Council of Finland [grant numbers 321716, 334814]. M.T.L. was supported by the Danish National Research Foundation [grant DNRF143].  L.L. was supported by Research Council of Finland [grant number 330887].

\printbibliography

\end{document}